% !TeX program = xelatex
\documentclass[UTF8,11pt]{ctexart}
\usepackage[a4paper,margin=2.15cm]{geometry}
\usepackage{amsmath,amssymb,mathtools}
\usepackage{mathrsfs}
\usepackage{booktabs,longtable,tabularx,array}
\usepackage{enumitem}
\usepackage{xcolor}
\usepackage[hidelinks]{hyperref}
\usepackage{url}
\usepackage{microtype}

\newcommand{\E}{\mathbb E}
\newcommand{\Pp}{\mathbb P}
\newcommand{\R}{\mathbb R}
\newcommand{\dd}{\,\mathrm d}
\newcommand{\sgn}{\operatorname{sgn}}
\newcommand{\supp}{\operatorname{supp}}
\newcommand{\Acat}{\textcolor{green!45!black}{\textbf{A：正确}}}
\newcommand{\Bcat}{\textcolor{blue!65!black}{\textbf{B：可由中文原稿完整修正}}}
\newcommand{\Ccat}{\textcolor{orange!80!black}{\textbf{C：须补充定义或条件修正}}}
\newcommand{\Dcat}{\textcolor{red!70!black}{\textbf{D：现有框架中难以修正}}}
\newcommand{\srcen}[1]{英文稿第 #1 行}
\newcommand{\srczh}[1]{中文稿 #1}
\setlist{nosep,leftmargin=2em}
\renewcommand{\arraystretch}{1.18}
\setlength{\tabcolsep}{4pt}

\title{\texttt{meanreversion\_article\_main.tex} 理论修正对比报告}
\author{逐定义、逐命题、逐模型的可复核数学审计}
\date{2026年8月14日}

\begin{document}
\maketitle

\begin{abstract}
本报告把两个问题严格分开：第一，英文稿是否忠实、完整地翻译了中文原稿；第二，中文原稿与英文稿中的数学陈述本身是否成立。审计对象是
\path{../meanreversion_article_main.tex}，对照源是工作区只读目录中的中文学位论文源文件
\path{chap/chap1.tex}、\path{chap/chap2.tex}、\path{chap/chap3.tex} 与
\path{chap/deno.tex}。结论不是简单的“翻译正确/错误”：英文稿基本逐段翻译了中文第三章正文，却删去了全部证明与依赖的基础定义；而中文证明本身又包含方向误用、局部鞅与真鞅混淆、定义域不相容和一个错误的 CIR 变量变换。

修正后可以保留的主线是：以双随机 Markov 核刻画一致 Copula；用有限时距条件期望恒等式作为严格的一阶均值回复准则；在扩展生成元域与真鞅条件下使用局部 PDE；用单调时变变换构造观测密度；再以独立平稳重置核构造混合 Copula。OU、CIR、仿射变换、指数变换和 Poisson 混合核在补齐条件后成立。Gamma 两例原先只验证一个零点而非全域符号；二次 CIR 仅能在标准 $\kappa>0$ 下作为有限时域单调模型；OU 特殊函数解在全实轴上必为非单射；CIR 的 Kummer 分离解需要指定分支、边界与可积性；其后的第二个变量变换则是错误的。实证注册表中的 25 个模型也不全属于修正后的一维 Markov 均值回复框架，必须区分直接嵌入、附条件嵌入、扩展状态嵌入和仅作诊断/基准四类。
\end{abstract}

\section{判定口径与最重要结论}

\subsection{四级分类}

本报告对每个数学对象给出唯一的主分类；“公式主体正确但原命题缺条件”归入 C，而不是 A。
\begin{description}[style=nextline]
  \item[\Acat] 在其原文已经明示的参数域和通常的条件期望约定下，陈述、公式与逻辑均成立；只需要排版或记号清理。
  \item[\Bcat] 英文稿有遗漏或翻译错误，但中文原稿已经给出足以得到正确版本的定义、参数域或证明；无需引入中文原稿之外的新数学条件。
  \item[\Ccat] 中文原稿也不充分或有逻辑漏洞；保留原研究思想，但必须改用更严谨的定义、补充支持集/定义域/可积性/边界/真鞅条件，或把“充分”降为“必要”。
  \item[\Dcat] 该对象不能在“实值一维 Markov 过程的全局严格递增变换”这一现有框架内成立；若要保留，必须改换模型类别（例如多分支隐状态或高维潜在状态），或删除错误推导。
\end{description}

\subsection{结论摘要}

\begin{enumerate}[label=\textbf{K\arabic*.}]
  \item 英文稿第 691 行开始的 \emph{Proofs} 为空；正文却引用当前文件中不存在的
  \texttt{copula-generator}、\texttt{time change copula}、\texttt{time change generator}、
  \texttt{dynamic martingale condition}、\texttt{E1} 与 \texttt{diffusion copula}。所以当前文件并非自包含的理论稿。
  \item 中文稿第三章第 55 行已经以注释“这里存疑”标记原均值回复漂移定义。该定义用
  $\E[X_t\mid X_{t-h}=x]$，后文却使用 $\E[X_{t+h}\mid X_t=x]$。修正版统一采用\emph{前向漂移}。
  \item “局部生成元等式为零”至多先给出局部鞅，不能无条件推出真鞅。中文稿第二章第 1200--1215 行及第三章第 747--759 行的“当且仅当”均需 class--(D)、一致可积或等价的 Dynkin 条件。
  \item 中文稿第三章第 706--710 行所谓命题证明把均值回复漂移直接写成零，实际复制了鞅证明，不能证明符号条件。
  \item 标准 OU/CIR 的条件均值、Gaussian-Copula 生成元、显式仿射/指数系数解、混合核和非局部债券 PDE 的公式主体成立；主要问题是条件和定义域没有写全。
  \item Gamma 两例只在目标均值对应的一个分位点令漂移为零。一个零点是全域符号条件的必要条件，不是充分条件。
  \item 二次 CIR 的 ODE 解正确，但当标准 CIR 的 $\kappa>0$、$A_0>0$ 时，线性系数最终必为负，故只能在满足显式单调约束的有限窗口内使用。
  \item OU 特殊函数候选依赖 $x^2$，在 OU 的状态空间 $\R$ 上是偶函数，除常数外不可能单射；CIR 的第二变量变换反代残差非零。这两项属于 D。
  \item 混合 Copula 的理论优势可以严格写成
  \[
    \E[Y_t-m\mid Y_s]
      =a_{s,t}R_{s,t}(Y_s-m),
    \qquad
    a_{s,t}=\exp\!\left\{-\int_s^t\lambda(u)\dd u\right\},
  \]
  即在完全保留一维边缘分布的同时，把条件均值偏离的衰减率从 $R_{s,t}$ 加快为 $a_{s,t}R_{s,t}$。这是清晰、可检验的理论优势；它不自动保证样本外密度分数优于所有基准。
\end{enumerate}

\section{原文映射与翻译完整性}

英文稿理论正文第 154--674 行基本对应中文 \path{chap/chap3.tex} 第 32--532 行。下表列出决定可复核性的遗漏。

\begin{longtable}{p{0.20\textwidth}p{0.21\textwidth}p{0.13\textwidth}p{0.38\textwidth}}
\toprule
对象 & 中文位置 & 分类 & 对英文稿的影响\\
\midrule
\endfirsthead
\toprule
对象 & 中文位置 & 分类 & 对英文稿的影响\\
\midrule
\endhead
Copula、Sklar 与广义分位数 & chap1:53--74；deno:5 & \Bcat & 可恢复基础定义；须另外说明不连续边缘下 Copula 不唯一。\\
$*$ 运算与一致族 & chap1:130--152 & \Ccat & 星积公式正确，但偏导应按 Lebesgue 几乎处处解释；仅有二元星一致性不足以判定一个既有过程的全部 Markov 性，宜直接以满足 Chapman--Kolmogorov 的双随机核为原始对象。\\
Copula 过程存在性 & chap1:154--178 & \Ccat & Dini 导数版本不自包含，$y=1$ 处还出现空下确界；修正版不依赖该写法。\\
OU/CIR 定义及 $\kappa>0$ & chap1:513--565 & \Bcat & 英文变换小节没有重新给出基础 SDE 和完整参数域；中文明确规定 $\kappa,\bar r,\sigma>0$，并给出 Feller 条件。\\
精确有限时距鞅条件 & chap2:51--56，1161--1181 & \Ccat & 核心恒等式正确，但只能对 Lebesgue 几乎处处的 $u\in(0,1)$ 断言；端点分位数通常无穷。\\
Copula 生成元 & chap2:60--75 & \Ccat & 可恢复右生成元思想；不能要求无界 Normal/Gamma 分位数属于 $C_0([0,1])$，也不应要求核一定有密度。\\
动态鞅 PDE & chap2:112--118，1183--1215 & \Ccat & 公式形式正确；证明分母有笔误，且“零局部漂移推出真鞅”缺 class--(D)/UI 条件。\\
Brownian/time-change 生成元 & chap2:125--140，1217--1283 & \Ccat & 公式正确；仅对 $t>0,h(t)>0$，不能声称精确定义域等于 $C_0^2$。\\
扩散 Copula 生成元 & chap2:154--162，1284--1309 & \Ccat & 公式主体正确；需内部光滑、零边界概率流。中文的全局一致椭圆条件排除 CIR，不能原样用于 CIR。\\
第三章全部证明 & chap3:704--883 & 混合 & 英文完全删除。平稳符号计算、强均值回复和债券生成元的核心可恢复；均值回复生成元证明、变换“iff”和混合命题证明需重写。\\
Kummer 函数附录 & chap3:887--906 & \Bcat & 英文只写 $\Phi,\Psi$ 而未定义，且与正态 CDF $\Phi$ 冲突；可由中文恢复后改记为 $M,U$，但模型分支有效性仍属 C。\\
CIR 第二变量变换 & chap3:439--452 & \Dcat & 中英文同错，不能通过忠实翻译修复；必须删除并换成正确降维。\\
\bottomrule
\end{longtable}

\section{符号逐项审计}

\begin{longtable}{p{0.13\textwidth}p{0.25\textwidth}p{0.13\textwidth}p{0.39\textwidth}}
\toprule
原符号 & 原用途/冲突 & 分类 & 修正版\\
\midrule
\endfirsthead
\toprule
原符号 & 原用途/冲突 & 分类 & 修正版\\
\midrule
\endhead
$\theta$ & 同时表示均值回复目标、OU/CIR 长期水平、Gamma 尺度及其极限 & \Ccat & 分别使用 $m$（目标均值）、$\bar r$（基础扩散长期水平）、$s$（Gamma 尺度）。原 CIR 显式式仅是 $m=\bar r$ 的特例。\\
$\alpha$ & Gamma 形状、负的回复系数；第 382 行又误写成正数 & \Ccat & Gamma 形状用 $\nu>0$；回复函数用 $\gamma(t)<0$，或 $d(t)=-\gamma(t)>0$。\\
$k$ & OU/CIR 速度、Kummer 函数的第一参数 & \Bcat & 回复速度统一用 $\kappa>0$；Kummer 参数用 $p,q$。\\
$c$ & Copula 密度、重置强度 $c(t)$、分离常数 & \Ccat & 分别用 $c_{s,t}(u,v)$、$\lambda(t)$、$\zeta$。\\
$A_t$ 与 $A(t)$ & Copula 生成元和变换系数视觉冲突 & \Bcat & 生成元写 $\mathcal A_t$；变换系数写 $a_1(t),a_2(t)$。\\
$\mu(t,x)$ & 被称为“reversion rate”，实际有状态量/时间量纲，是局部漂移 & \Ccat & 写作 $b_X(t,x)$，称前向局部漂移；线性回复速度另记 $d(t)>0$。\\
$\operatorname{Ran}(X_t)$ & 随机变量的“值域”依版本而变，且会包括零概率点 & \Ccat & 使用边缘分布的拓扑支持 $E_t=\supp(F_t)$；等式按核版本对 $F_t$--a.e. 状态成立，连续时再延拓。\\
$F_t^{-1}$ & 端点处可能为 $\pm\infty$，却要求属于 $C_0^2([0,1])$ & \Ccat & 明确定义广义分位数 $Q_t$；在 $(0,1)$ 上使用扩展/加权生成元域，并单列端点可积性。\\
$f_t$ 与 $f(t,x)$ & 前者是边缘密度，后者是时变变换 & \Bcat & 密度用 $p_t$ 或 $q_t$；变换保留 $f$。\\
$\Phi$ & 同时是正态 CDF 和 Kummer 第一解 & \Bcat & 正态 CDF 保留 $\Phi$；Kummer 解写 $M(p,q,z)={}_1F_1(p;q;z)$ 与 $U(p,q,z)$。\\
$Y$ & 既是过程又是债券 PDE 中独立平稳抽样变量 & \Bcat & 过程用 $Y_t$；积分变量的法律写 $\pi(\dd z)$，不再引入同名随机变量。\\
$\mathbb R^+$ & 是否包含零不明确，CIR 边界尤其敏感 & \Ccat & 时间写 $[0,\infty)$；正实数写 $(0,\infty)$；CIR 状态写 $[0,\infty)$。\\
条件 $X_t=x$ & 连续变量事件通常概率为零 & \Ccat & 全文均解释为选定的正则条件转移核；结论先按边缘分布几乎处处成立。\\
\bottomrule
\end{longtable}

\section{英文稿全部具名对象与带标签公式的逐项判定}

下表覆盖英文稿理论正文第 154--674 行的全部 \verb|\label{...}|。位置栏中的行号以原英文文件为准；“正确部分”只表示该局部公式可保留，不覆盖同一段落的过强文字结论。

\small
\begin{longtable}{p{0.055\textwidth}p{0.20\textwidth}p{0.095\textwidth}p{0.14\textwidth}p{0.42\textwidth}}
\toprule
序号 & 原标签/对象 & 原行 & 分类 & 严格结论与修正动作\\
\midrule
\endfirsthead
\toprule
序号 & 原标签/对象 & 原行 & 分类 & 严格结论与修正动作\\
\midrule
\endhead
1 & \texttt{mean-E2} & 158--160 & \Acat & OU 与标准 CIR 在 $\kappa>0$ 时均满足该条件均值式；CIR 还要有非负解及一阶矩。\\
2 & \texttt{mean-E3} & 162--165 & \Acat & 由上一式直接得到；从条件均值到无条件均值须 $X_s$ 可积。\\
3 & Definition \texttt{mean-reversion-def} & 170--188 & \Ccat & 目标思想可保留；把 backward drift 改为基于 $P_{t,t+h}$ 的 forward drift，状态用支持集，条件等式按正则条件核的版本解释。\\
4 & \texttt{mean-reversion-\allowbreak def-cond-1} & 174--176 & \Acat & 无条件一阶矩收敛是清楚且可保留的条件；但它不蕴含长期条件均值收敛。\\
5 & \texttt{drift-condition} & 179--181 & \Ccat & $t=0$ 无定义，且与后文生成元方向不一致。改为 $\lim_{h\downarrow0}\{\int yP_{t,t+h}(x,\dd y)-x\}/h$。\\
6 & \texttt{mean-reversion-\allowbreak condition-1} & 183--186 & \Ccat & 符号思想正确；应写 $(x-m)b_X(t,x)<0$（$x\ne m$）及 $b_X(t,m)=0$，并明确几乎处处/连续版本。\\
7 & Assumption \texttt{hyp1} & 205--211 & \Ccat & 对全 $\R$ 严格递增排除了 Gamma/CIR；连续 Copula 密度又排除奇异核。改为支持内部严格递增，并以双随机 Markov 核为主。\\
8 & Assumption \texttt{hyp2} & 212--221 & \Ccat & Normal/Gamma 分位数在端点无界，不在原 $C_0$ 域。使用 time--space 扩展生成元、局部化和可积包络。\\
9 & \texttt{hyp2 relation} & 217--219 & \Ccat & 这是人为拼接左右极限的强条件；修正版直接用前向生成元，不再需要它。\\
10 & Proposition \texttt{mean-reversion-\allowbreak def-1} & 222--235 & \Ccat & 结论可在前向漂移恒等式成立时修复。中文证明第 709 行把漂移误写成零；正确式为 $b_X(t,Q_t(u))=\partial_tQ_t+\mathcal A_tQ_t$。\\
11 & \texttt{mean-reversion-\allowbreak condition-2} & 231--232 & \Ccat & 正确的生成元符号形式，但只在扩展生成元域、交换极限合法及 $u\in(0,1)$ 几乎处处时成立。\\
12 & Example \texttt{gamma mean reversion example} & 243--314 & \Ccat & 两个 Case 都只能保留为候选族与必要条件诊断，不能宣称已证明均值回复。\\
13 & \texttt{OU generator} & 247--249 & \Acat & 公式等于 $a\{g''(z)-zg'(z)\}$ 的分位坐标形式，代数正确；定义域由第 7--9 项修正。\\
14 & \texttt{gamma example 1} & 255--260 & \Ccat & 该标量方程只保证目标均值分位点的漂移为零；还需验证所有 $u\in(0,1)$ 的符号及端点可积性。\\
15 & \texttt{gamma example 2} & 292--295 & \Ccat & ODE 仍只匹配移动目标分位点；应降为必要条件，并另加全域符号、尺度正性、极限和 ODE 解存在性。\\
16 & \texttt{gamma example 3} & 298--300 & \Acat & 若全域符号成立，则目标点漂移为零确为必要条件；原稿错误在把此必要条件反向当成充分条件。\\
17 & Proposition \texttt{stationary distribution m.r.} & 317--321 & \Ccat & $g_F''-zg_F'$ 的全域符号是正确充分条件；须删除 $F^{-1}\in C_0^2([0,1])$，改为内部扩展域与端点可积/局部化。\\
18 & Proposition \texttt{quantile-equation} & 328--334 & \Ccat & 核心正确；补 $X_t$ 可积、条件期望相对于 $\mathcal F_s$、$R_{s,t}>0$ 乘法性和极限交换。\\
19 & \texttt{mean reversion condition-1} & 330--333 & \Acat & 这是严格且最有用的有限时距定义；$\gamma<0$ 且 $R\to0$ 时同时推出局部方向和长期条件均值。\\
20 & \texttt{XY relation} & 341--343 & \Acat & 确定性非零仿射变换正确，且保留 Markov 性；真鞅仍需可积。\\
21 & \texttt{y martingale equation} & 347--349 & \Ccat & 是正确局部 PDE；无 class--(D)/UI 时最多得到局部鞅，不能写无条件 iff。\\
22 & \texttt{XY relation-1} & 351--353 & \Acat & 广义分位数在递增仿射变换下的关系正确。\\
23 & \texttt{mean reversion \allowbreak condition-2} & 355--357 & \Ccat & 代数正确；其全局充分性继承第 21 项的真鞅条件。\\
24 & \texttt{mean reversion \allowbreak condition-4} & 364--367 & \Acat & 这是强均值回复的精确 Copula--分位数条件；密度不存在时把 $c_{s,t}(u,v)\dd v$ 改为核 $K_{s,t}(u,\dd v)$。\\
25 & \texttt{diffusion process} & 374--376 & \Ccat & 通用 SDE 还缺存在唯一性与状态区间条件；中文的一致椭圆条件不能覆盖 CIR，须把 CIR 单列。\\
26 & Proposition \texttt{marginal mean reversion} & 392--397 & \Ccat & It\^o 推导给出 PDE 是正确的；“当且仅当”要补严格递增、支持映射、真鞅和矩条件。\\
27 & \texttt{mean reversion \allowbreak condition-3} & 394--396 & \Acat & 写成 $(\partial_t+\mathcal L)f=\gamma(t)(f-m)$ 后符号正确；第 382--384 行把 $\gamma<0$ 误写为正数，应改正。\\
28 & \texttt{OU copula equation-1} & 403--405 & \Ccat & 当基础 OU 与目标共享同一个 $\theta$ 时代数成立；修正版分离 $m$ 和基础水平，或先中心化基础 OU。\\
29 & Example \texttt{OU example} & 407--477 & 混合 & 仿射、指数项为 A/C；特殊函数项作为 PDE 解为 A，但作为全轴单调模型为 D，不能整体宣称三者均有效。\\
30 & \texttt{OU copula equation} & 409--411 & \Acat & 在基础 OU 中心为零、目标 $m=0$、$\gamma=-d$ 时正确。\\
31 & \texttt{OU empirical} & 427--432 & \Acat & $A_0>0$ 时全轴严格递增，正态边缘均值/方差正确；$\kappa>d$ 时方差增长并不违反一阶条件均值回复。\\
32 & \texttt{OU example-2 parameter} & 447--453 & \Acat & 三个 ODE 解的反代残差为零。\\
33 & \texttt{OU example-2} & 455--465 & \Ccat & $A_0>0$ 给出单调性；正支持还需 $m+C(t)\ge0$。密度支持为 $y>m+C(t)$，并需写 Jacobian。\\
34 & OU 特殊函数（无单独标签） & 467--475 & \Dcat & $f(t,x)$ 只依赖 $x^2$，故 $f(t,x)=f(t,-x)$；不能在 OU 的 $\R$ 状态空间上作为非平凡单调 Copula 保持变换。\\
35 & \texttt{CIR-pde-1} & 483--485 & \Ccat & 原文把目标和 CIR 长期水平都记作 $\theta$；修正版分别写 $m,\bar r$。\\
36 & Example \texttt{CIR empircial example} & 487--598 & 混合 & 四组候选的有效层级不同，必须拆开判定，不能用一个例子统一背书。\\
37 & \texttt{CIR-pde} & 489--491 & \Acat & 在 $m=\bar r=\theta$、$\gamma=-d$ 的特例下代数正确。\\
38 & \texttt{CIR empirical} & 511--514 & \Acat & 仿射系数解正确；$A_0>0$ 时在 $[0,\infty)$ 单调，支持和基础 CIR 参数域必须明示。\\
39 & \texttt{CIR example-2} & 517--522 & \Acat & 二次 ansatz 的三项 ODE 正确。\\
40 & \texttt{CIR example-2 parameter} & 533--541 & \Acat & 在 $\kappa\ne0$ 时显式系数反代残差均为零；这不等于全时域单调。\\
41 & \texttt{CIR empirical-2} & 543--546 & \Ccat & 标准 $\kappa>0$ 下只能在有限 $[0,T]$ 使用，并要求 $B_0\ge\eta(e^{\kappa T}-1)$；严格正 Jacobian 宜取严格不等式。\\
42 & \texttt{CIR example-3 parameter} & 566--574 & \Ccat & Riccati 解正确；须有 $A_0>0$、分母不为零，平稳 Gamma 启动时还需 $A_0<2\kappa/\sigma^2$。\\
43 & \texttt{CIR empirical-3} & 576--579 & \Ccat & 单调分支可构造；实际支持为 $y\ge m+C(t)+e^{B(t)}$，而不是只写 $y>C(t)$，并需指数矩/真鞅条件。\\
44 & CIR Kummer 分离解（无单独标签） & 581--583 & \Ccat & 分离方程正确；须固定 $M/U$ 分支、常数、边界、初始分布、单调性和可积性。$M(p,q,z)$ 的 $p,q>0$ 分支虽递增，却对平稳 Gamma 不可积。\\
45 & CIR 第二变量变换（无单独标签） & 584--596 & \Dcat & 反代残差一般非零；后列三组 $u$ 不能据此视为 CIR PDE 解。应删除并以含一阶项的正确降维替换。\\
46 & 统一结论段 & 600 & \Ccat & 仿射项可全时域成立；二次项只在有限窗口；指数项需矩条件。原来的统一“均为单调、均值趋于 $\theta$”过强。\\
47 & \texttt{mix copula expression} & 605--608 & \Acat & 在基础核保持同一不变分布且 $a_{s,t}$ 乘法时，混合核满足 Chapman--Kolmogorov。\\
48 & \texttt{mix-copula-process} & 611--614 & \Ccat & 构造正确，但“由一致性即存在”需改为先给定双随机 Markov 核；可包含无密度的奇异分量。\\
49 & Proposition \texttt{mix copula mean reversion} & 618--620 & \Ccat & 结论可直接由混合转移核证明。中文证明错误地反用一个充分命题，并在最后把 $u$ 与物理量 $\theta$ 比较。\\
50 & Proposition \texttt{mix mean reversion construction} & 629--635 & \Ccat & Poisson 构造及最终核正确；中文第 802 行重复乘事件概率，证明须重写。非齐次强度也可严格扩展。\\
51 & \texttt{mean reversion construction} & 631--633 & \Acat & 以 $N_t$ 选择独立平稳副本的定义有效；独立、同分布、平稳启动缺一不可。\\
52 & Example \texttt{mix copula construction} & 649--655 & \Acat & 从平稳 OU 副本和独立 Poisson 时钟得到混合 OU 核；参数要求 $\kappa,\sigma,\lambda>0$。\\
53 & \texttt{mean reversion OU construction} & 651--653 & \Acat & 与第 51 项相同的正确具体化。\\
54 & Proposition \texttt{mix copula bond pde} & 659--672 & \Ccat & 非局部生成元和 PDE 符号正确；须说明风险中性测度、指数泛函可积、解类别和唯一性。\\
55 & \texttt{bond equation definition} & 661--663 & \Acat & 风险中性条件期望定义正确；需把测度明确记为 $\mathbb Q$。\\
56 & \texttt{mix copula bond price pde} & 665--670 & \Ccat & 核心为 $g_t+\mathcal Gg-yg=0$，正确；积分项宜写 $\int g(t,z)\pi(\dd z)$，避免同名随机变量。\\
\bottomrule
\end{longtable}
\normalsize

\paragraph{完整性说明。}
上表 1--56 已包含原理论区间的全部 52 个带标签对象，以及四个没有单独标签但承载独立理论结论的段落。空的 Introduction、Empirical Analysis、Conclusion 和 Proofs 不属于可判定的数学对象；“Proofs 为空”已在翻译完整性章节单列。

\section{关键定义与定理为什么必须重写}

\subsection{一致 Copula 应先定义为双随机 Markov 核}

令 $I=(0,1)$，$\ell$ 为其 Lebesgue 测度。一族核 $\{K_{s,t}:0\le s\le t\}$ 若满足
\begin{align}
K_{t,t}(u,\cdot)&=\delta_u,\label{cmp:kid}\\
K_{s,t}(u,A)&=\int_I K_{r,t}(v,A)K_{s,r}(u,\dd v),\qquad s<r<t,\label{cmp:kck}\\
\int_I K_{s,t}(u,A)\ell(\dd u)&=\ell(A),\label{cmp:bistochastic}
\end{align}
则它是双随机 Markov 演化；由
\[
C_{s,t}(x,y)=\int_0^x K_{s,t}(u,(0,y])\dd u
\]
定义的二元分布是 Copula。这个定义同时覆盖连续密度、确定性核和 Poisson “无重置”产生的奇异分量。若密度存在，才写
$K_{s,t}(u,\dd v)=c_{s,t}(u,v)\dd v$。

中文稿的 $*$ 公式在绝对连续情形下可恢复为
\[
(C*D)(u,v)=\int_0^1 \partial_2C(u,w)\,\partial_1D(w,v)\dd w,
\]
其中偏导仅需几乎处处存在。仅陈述所有\emph{二元} Copula 的星一致性不足以判定一个已经给定、可能具有不同高维联合分布的过程本身为 Markov；安全陈述是：满足 \eqref{cmp:kck} 的核构造一个 Markov 实现。

\paragraph{英文修订稿采用的简化范围。}
上述核表述给出一般理论边界；但本文英文理论及其实证似然只需要正时距
$s<t$ 的绝对连续 Copula。因此更新稿进一步假设
$K_{s,t}(u,\dd v)=c_{s,t}(u,v)\dd v$，并把双随机性与
Chapman--Kolmogorov 直接写成密度的行列归一化和卷积恒等式。对角
$s=t$ 仍是 Dirac 恒等算子，不要求具有密度。该范围选择排除未单独处理的
奇异 Copula，但不排除本文的 OU/CIR 独立重置模型，因为其正时距密度为
$c^{\rm mix}_{s,t}=a_{s,t}c_{s,t}+(1-a_{s,t})$。

\subsection{有限时距恒等式优先于局部 PDE}

设 $Q_t=F_t^{-1}$，$U_t=F_t(X_t)$。对 $X_t$ 可积的 Markov 过程，最稳健的真鞅条件是
\begin{equation}
Q_s(u)=\int_I Q_t(v)K_{s,t}(u,\dd v)
\quad\text{for }F_s\text{-a.e. }u,\label{cmp:exact-mg}
\end{equation}
而强一阶均值回复的精确条件是
\begin{equation}
\int_I\{Q_t(v)-m\}K_{s,t}(u,\dd v)
=R_{s,t}\{Q_s(u)-m\}.\label{cmp:exact-mr}
\end{equation}
这两个式子本身已经是有限时距条件期望，不需要从局部鞅再推回真鞅。

只有当 $Q(t,u)$ 属于 time--space 扩展生成元域、极限可交换时，才有
\begin{equation}
b_X(t,Q_t(u))=(\partial_t+\mathcal A_t)Q(t,u).\label{cmp:forward-generator}
\end{equation}
若 $(\partial_t+\mathcal A_t)Q=0$，一般先得到局部鞅；要推出真鞅，至少需要每个有限 $T$ 上
$\{Q(\tau,U_\tau):\tau\le T\}$ 为 class--(D)，或给出足以应用 Dynkin 公式和一致可积性的增长条件。中文证明第二章第 1187、1192 行还把前向差商分母误写为 $t-s$，正确分母是 $T-t$。

\subsection{扩散变换命题的严格版本}

设基础扩散
\[
\dd r_t=b(r_t)\dd t+a(r_t)\dd W_t,
\]
令 $Y_t=f(t,r_t)$，且 $f(t,\cdot)$ 在基础状态区间上严格递增。若
\begin{equation}
(\partial_t+\mathcal L)f(t,x)=\gamma(t)\{f(t,x)-m\},
\qquad \gamma(t)<0,\label{cmp:transform-pde}
\end{equation}
并且对每个 $T<\infty$
\begin{equation}
\E\int_0^T R_{0,t}^{-2}a(r_t)^2f_x(t,r_t)^2\dd t<\infty,
\qquad
R_{s,t}=\exp\!\left\{\int_s^t\gamma(u)\dd u\right\},\label{cmp:square-int}
\end{equation}
则 It\^o 公式给出
\[
\E[Y_t-m\mid\mathcal F_s]=R_{s,t}(Y_s-m).
\]
反向结论要用半鞅漂移分解的唯一性，并要求同样的正则性。原文只写 $f\in C^{1,2}$，不足以保证局部随机积分是真鞅。

若基础转移密度是 $p_{s,t}(x_s,x_t)$，则观测密度为
\begin{equation}
q_{s,t}(y_s,y_t)=p_{s,t}\!\left(f_s^{-1}(y_s),f_t^{-1}(y_t)\right)
\left|\frac{\partial f_t^{-1}(y_t)}{\partial y_t}\right|.\label{cmp:jacobian}
\end{equation}
这要求逆像唯一、观测落在 $f_t(E)$ 内，并且 Jacobian 在内部有限；“PDE 解存在”本身并不等于“观测密度可构造”。

\section{Gaussian--Gamma 两例的逻辑缺口}

令 $G_\nu$ 为 $\Gamma(\nu,1)$ 的分位数，$s>0$ 为尺度，
$Q(u)=sG_\nu(u)$，并记
\[
u_\nu=F_{\nu,1}(\nu),\qquad z_\nu=\Phi^{-1}(u_\nu).
\]
原 Case 1 推出的标量条件
\begin{equation}
\Gamma(\nu)\left(\frac e\nu\right)^\nu\phi(z_\nu)-2z_\nu=0
\label{cmp:gamma-root}
\end{equation}
代数正确；它等价于只在 $Q(u_\nu)=\nu s$ 处要求 $\mathcal A Q=0$。严格均值回复却要求
\begin{equation}
\sgn\{\mathcal A Q(u)\}=\sgn\{\nu s-Q(u)\}
\quad\text{for every }u\in(0,1)\text{（至少 a.e.）}.\label{cmp:gamma-global}
\end{equation}
因此 \eqref{cmp:gamma-root} 只是 \eqref{cmp:gamma-global} 的必要条件。原文没有证明零点唯一，也没有证明两侧符号。已有高精度数值扫描未发现有限的 $\nu$ 使 \eqref{cmp:gamma-root} 精确为零，但这不能替代解析不存在性证明；正确做法是把该族保留为“全域符号诊断失败/未证”的候选，而不是声称存在均值回复模型。

对时变尺度 $Q_t(u)=s(t)G_\nu(u)$，定义
\[
\mathscr G_\nu(u)=2\phi(\Phi^{-1}u)\Phi^{-1}(u)G_\nu'(u)
-\phi(\Phi^{-1}u)^2G_\nu''(u).
\]
当 Gaussian Copula 的速率为 $a>0$ 时，真实前向漂移是
\begin{equation}
D_t(u)=s'(t)G_\nu(u)-a s(t)\mathscr G_\nu(u).\label{cmp:gamma-tv-drift}
\end{equation}
若 $s(t)\to s_\infty$，目标均值应是 $m=\nu s_\infty$，全域必要充分符号条件为
\begin{equation}
\sgn D_t(u)=\sgn\{m-s(t)G_\nu(u)\}.\label{cmp:gamma-tv-global}
\end{equation}
原 ODE 只是把 \eqref{cmp:gamma-tv-drift} 在满足 $s(t)G_\nu(u_t)=m$ 的一个 $u_t$ 上设为零。即使 ODE 有正解并收敛，极限端点仍须满足 Case 1 的零点条件，故 Case 2 不能绕开 Case 1 的可行性问题。

\section{OU 与 CIR 变换逐模型核验}

\subsection{基础过程与参数域}

中心化 OU 写为
\[
\dd r_t=-\kappa r_t\dd t+\sigma\dd W_t,\qquad \kappa>0,\ \sigma>0.
\]
其精确转移是 Normal。一般均值为 $\bar r$ 的 OU 只需把 $r_t$ 换成 $r_t-\bar r$。标准 CIR 写为
\begin{equation}
\dd r_t=\kappa(\bar r-r_t)\dd t+\sigma\sqrt{r_t}\dd W_t,
\qquad \kappa,\bar r,\sigma>0,\quad r_t\in[0,\infty).
\label{cmp:cir-sde}
\end{equation}
条件转移可写为
\[
r_t\mid r_s=x\ \overset d=
\frac{\sigma^2(1-e^{-\kappa\Delta})}{4\kappa}
\chi^{\prime 2}_{\,4\kappa\bar r/\sigma^2}
\!\left(\frac{4\kappa e^{-\kappa\Delta}x}
{\sigma^2(1-e^{-\kappa\Delta})}\right),\qquad \Delta=t-s.
\]
$2\kappa\bar r\ge\sigma^2$ 是零边界不可达的简洁充分条件；CIR 非负解和内部转移密度在更弱参数下仍可讨论，但边界处理必须单列。中文稿把这些条件放在第一章，英文变换小节未移植。

以下令 $d>0$，目标均值为 $m$，并使用 PDE
\[
(\partial_t+\mathcal L)f+d(f-m)=0.
\]

\subsection{OU：两类有效解和一个结构失败}

\paragraph{仿射 OU。}
\[
f(t,x)=m+A_0e^{(\kappa-d)t}x+B_0e^{-dt},\qquad A_0>0.
\]
它在全实轴严格递增，ODE 残差为零，且多项式增长保证有限时域真鞅条件。若基础 OU 平稳，原稿给出的 Normal 均值和方差正确。$\kappa>d$ 时方差增长，说明一阶条件均值回复不蕴含平稳或方差有界，而不是公式错误。

\paragraph{指数 OU。}
\begin{align*}
f(t,x)&=m+\exp\{A(t)x+B(t)\}+C(t),\\
A(t)&=A_0e^{\kappa t},\\
B(t)&=B_0-\frac{\sigma^2A_0^2}{4\kappa}(e^{2\kappa t}-1)-dt,\\
C(t)&=C_0e^{-dt}.
\end{align*}
$A_0>0$ 给出严格递增，Normal 指数矩给出每个有限时域的可积性；支持为
$y>m+C(t)$。若要把它作为正价格模型，还需 $m+C(t)\ge0$，原稿仅写 $A(t)>0$ 不够。

\paragraph{OU Kummer 候选。}
\[
f(t,x)=m+e^{-\zeta t}
J\!\left(\frac{\zeta-d}{2\kappa},\frac12;
\frac{\kappa x^2}{\sigma^2}\right)
\]
确实把 PDE 化成 Kummer 方程，故是\emph{代数上的 PDE 解}。但它关于 $x$ 为偶函数；任意非恒定分支都在 $\R$ 上非单射。用两逆像密度可以定义另一个观测模型，却不再保持原 Copula，且观测过程一般不再是一维 Markov。因此在现有理论框架中判为 D。

\subsection{CIR：仿射、二次、指数和特殊函数}

\paragraph{仿射 CIR。}
令
\begin{align*}
a_1(t)&=A_0e^{(\kappa-d)t},\\
b_1(t)&=m-A_0\bar r e^{(\kappa-d)t}+B_0e^{-dt},\\
f(t,x)&=a_1(t)x+b_1(t).
\end{align*}
这里 $B_0$ 是重参数化常数；当 $m=\bar r$ 时退化为原稿公式。$A_0>0$ 时全半轴严格递增，支持是 $[b_1(t),\infty)$，转移密度由 \eqref{cmp:jacobian} 精确得到。

\paragraph{二次 CIR。}
令
\[
f(t,x)=A(t)x^2+B(t)x+C(t),\qquad
\eta=\frac{A_0(\sigma^2+2\kappa\bar r)}{\kappa}.
\]
正确系数为
\begin{align*}
A(t)&=A_0e^{(2\kappa-d)t},\\
B(t)&=e^{(\kappa-d)t}\{B_0-\eta(e^{\kappa t}-1)\},\\
C(t)&=m+\frac{\bar r\eta}{2}e^{(2\kappa-d)t}
-\bar r(B_0+\eta)e^{(\kappa-d)t}
+\left(C_0-m+\bar rB_0+\frac{\bar r\eta}{2}\right)e^{-dt}.
\end{align*}
三条 ODE 残差均为零。因为 $f_x=2A(t)x+B(t)$，当标准参数
$\kappa,\bar r,\sigma,A_0>0$ 时，$A(t)>0$ 且 $\eta>0$；在固定窗口 $[0,T]$ 上单调非减的充要条件是
\begin{equation}
B_0\ge\eta(e^{\kappa T}-1).
\label{cmp:quadratic-window}
\end{equation}
但右端随 $T\to\infty$ 发散，任何有限 $B_0$ 都不能给出全时域单调性。若要保证支持内部 Jacobian 处处严格为正，应在 \eqref{cmp:quadratic-window} 取严格不等式。逆变换取唯一非负根
\[
x=\frac{2(y-C)}{B+\sqrt{B^2+4A(y-C)}},
\qquad \left|\frac{\partial x}{\partial y}\right|=\frac1{2Ax+B}.
\]

\paragraph{$\kappa<0$ 不能作为同一标准 CIR 的补救。}
纯代数上，当 $\kappa<0$ 时二次映射可在
\[
A_0>0,\qquad \min\{B_0,B_0+\eta\}\ge0
\]
下全时域单调。但这不满足原中文第一章明确规定的标准 CIR 参数域。若 $\bar r>0$，零点处漂移 $\kappa\bar r<0$ 会把状态推出非负半轴；若令 $\bar r<0$ 使截距为正，漂移则随状态向上发散，不存在原 Gamma 平稳分布或均值回复。因此 $\kappa<0$ 至多定义另一个超临界平方根扩散模型，不能把现有 CIR 实证的参数约束简单放宽。

\paragraph{指数 CIR。}
\begin{align*}
f(t,x)&=m+\exp\{A(t)x+B(t)\}+C(t),\\
A(t)&=\frac{2\kappa}{\sigma^2+(2\kappa/A_0-\sigma^2)e^{-\kappa t}},\\
B(t)&=-dt-\frac{2\kappa\bar r}{\sigma^2}
\log\!\left[(2\kappa/A_0-\sigma^2)+\sigma^2e^{\kappa t}\right]
+B_0+\frac{2\kappa\bar r}{\sigma^2}\log(2\kappa/A_0),\\
C(t)&=C_0e^{-dt}.
\end{align*}
这里把 $C(t)$ 定义为相对于目标 $m$ 的附加项；原文的总加法项为
$m+(C_0-m)e^{-dt}$，两种参数化等价。$A_0>0$ 保证有限时点 $A(t)>0$；若基础 CIR 从平稳 Gamma 分布启动，为保证指数矩有限还要
\begin{equation}
0<A_0<\frac{2\kappa}{\sigma^2}.
\label{cmp:cir-exp-moment}
\end{equation}
支持的精确下界是 $m+C(t)+e^{B(t)}$。原稿写成仅 $y>C(t)$，漏掉了 CIR 状态 $x\ge0$ 带来的 $e^{B(t)}$。

\paragraph{CIR Kummer 分离族。}
令
\[
p=\frac{\zeta-d}{\kappa},\qquad q=\frac{2\kappa\bar r}{\sigma^2},\qquad
z=\frac{2\kappa x}{\sigma^2}.
\]
则
\[
f(t,x)=m+e^{-\zeta t}J(p,q,z),\qquad
zJ_{zz}+(q-z)J_z-pJ=0
\]
的分离推导正确。可是 $J=C_1M(p,q,z)+C_2U(p,q,z)$ 没有指定分支、边界和归一化，并不定义一个模型。$p,q>0,C_1>0,C_2=0$ 时 $M$ 分支严格递增，但其 $e^z$ 尾部使其对平稳 $z\sim\Gamma(q,1)$ 不可积。其他分支必须逐一证明单调、支持和矩条件，不能用一个“特殊函数模型”统称。

\paragraph{错误的第二变换及正确替代。}
原稿声称某指数变量替换把方程化为 $u_\tau=zu_{zz}$，反代一般残差非零。一个直接、可验证的替代是
\[
f-m=e^{-dt}u(\tau,z),\qquad
z=xe^{\kappa t},\qquad
\tau=\frac{\sigma^2}{2\kappa}(1-e^{\kappa t}),
\]
它给出
\begin{equation}
u_\tau=zu_{zz}+\frac{2\kappa\bar r}{\sigma^2}u_z.
\label{cmp:correct-cir-reduction}
\end{equation}
一阶项一般不能消失。原稿列在错误变换之后的三个 $u$ 即使满足另一个 PDE，也没有因此成为原 CIR 方程的解。

\section{混合 Copula、Poisson 重置与债券 PDE}

令基础 Markov 核 $P_{s,t}$ 保持同一不变概率 $\pi$，并定义独立重置核
$\Pi(x,A)=\pi(A)$。对非负局部可积强度 $\lambda(t)$，令
\[
a_{s,t}=\exp\!\left\{-\int_s^t\lambda(u)\dd u\right\},\qquad
P^{\rm mix}_{s,t}=a_{s,t}P_{s,t}+(1-a_{s,t})\Pi.
\]
由于 $a_{s,r}a_{r,t}=a_{s,t}$、$P_{s,r}P_{r,t}=P_{s,t}$ 及
$P\Pi=\Pi P=\Pi$，直接相乘得到
\[
P^{\rm mix}_{s,r}P^{\rm mix}_{r,t}=P^{\rm mix}_{s,t}.
\]
所以一致性结论正确，并且不需要假定转移核有连续密度。其生成元是有界重置扰动
\begin{equation}
\mathcal G_t^{\rm mix}h(x)=\mathcal G_t h(x)
+\lambda(t)\left\{\int h(z)\pi(\dd z)-h(x)\right\}.
\label{cmp:mixed-generator}
\end{equation}

若基础过程满足强一阶均值回复
\[
\E[Z_t-m\mid Z_s=x]=R_{s,t}(x-m),
\]
则混合过程精确满足
\begin{equation}
\E[Y_t-m\mid Y_s=x]=a_{s,t}R_{s,t}(x-m).
\label{cmp:mixed-mr}
\end{equation}
若 $R_{s,t}=\exp\{\int_s^t\gamma(u)\dd u\}$，混合后的瞬时系数是
$\gamma(t)-\lambda(t)$，局部漂移为
\[
b_{\rm mix}(t,x)=b_0(t,x)+\lambda(t)(m-x).
\]
因此重置项严格加强偏离的条件均值衰减，同时保持所有一维边缘为 $\pi$。这就是混合 Copula 主张中最清晰的结构性优势。

原中文混合命题的结论可以这样直接证明，但其第 765--768 行把一个“符号条件推出均值回复”的充分命题反向使用；第 791 行还把分位数 $u\in[0,1]$ 与物理状态 $m$ 直接比较。Poisson 构造的最终核正确，但第 802 行把联合条件概率与 Poisson 权重重复相乘。修正版全部使用上面的核证明。

令 $N$ 为独立的非齐次 Poisson 过程，强度 $\lambda(t)$；令
$\{Z^{(j)}\}_{j\ge0}$ 是独立同分布的平稳基础过程，并设
$Y_t=Z_t^{(N_t)}$。区间 $(s,t]$ 无跳时沿原副本演化，至少一跳时终值来自独立平稳副本，故恰得到 $P^{\rm mix}_{s,t}$。常强度是原稿命题，非齐次版本严格容纳实证中的 SeasonalIntensityMOU。

在风险中性测度 $\mathbb Q$ 下，以平稳 OU 为基础时，
\[
\mathcal G h(y)=\kappa(m-y)h'(y)+\frac12\sigma^2h''(y)
+\lambda\left\{\int h(z)\pi(\dd z)-h(y)\right\}.
\]
若指数泛函可积且 Feynman--Kac 的经典解条件（或相应黏性解比较原则）成立，则
\[
g(t,y)=\E^{\mathbb Q}\!\left[
e^{-\int_t^T Y_u\dd u}\mid Y_t=y\right]
\]
满足
\begin{equation}
g_t+\kappa(m-y)g_y+\tfrac12\sigma^2g_{yy}
+\lambda\left\{\int g(t,z)\pi(\dd z)-g(t,y)\right\}-yg=0,
\qquad g(T,y)=1.
\label{cmp:bond-pde}
\end{equation}
原 PDE 的符号和非局部项正确；原命题只写 $C^{1,2}$，没有给风险中性测度、增长、可积性和唯一性，故归入 C。

\section{中文证明与基础材料中的可定位错误}

\begin{longtable}{p{0.17\textwidth}p{0.14\textwidth}p{0.60\textwidth}}
\toprule
位置 & 分类 & 错误及修正\\
\midrule
\endfirsthead
\toprule
位置 & 分类 & 错误及修正\\
\midrule
\endhead
chap1:119--152 & \Ccat & 连续边缘未先规定时，Sklar Copula 不唯一；二元星一致性与一个既有过程的全有限维 Markov 分解不能混为一谈。\\
chap1:156--178 & \Ccat & modified Dini 定义截掉边界分支；原文在 $y=1$ 处取 $\inf_{v>1}$，集合为空。修正版改用双随机核。\\
chap1:520 & \Bcat & OU 自然滤过误写为 $\sigma(X_u)$，应为 $\sigma(r_u)$。\\
chap1:544 & \Bcat & 零边界不可达条件可写 $2\kappa\bar r\ge\sigma^2$，不必严格大于。\\
chap1:553--555 & \Bcat & Gamma 混合密度的幂次少 $-1$；应为 $z^{i+\nu/2-1}$。\\
chap1:564--565 & \Ccat & 一般非平稳初值下 CIR Copula 密度把时点边缘逆 CDF 的非中心参数写成 0，不可原样移植。修正版直接给正确 CIR 转移密度。\\
chap2:39--48 & \Ccat & “后续全部模型满足 hyp1”错误：Gamma/CIR 支持不是全 $\R$；Poisson 无跳分量还会产生奇异核。\\
chap2:60--88 & \Ccat & $C_0$ 在符号表定义于 $\R$，却写成 $C_0([0,1])$；无界分位数不在其中。\\
chap2:112--118 & \Ccat & 动态鞅 PDE 形式正确，但充分性缺 true-martingale 条件。\\
chap2:1187,1192 & \Bcat & 前向差商分母 $t-s$ 是直接笔误，应为 $T-t$；即使修正仍需上一行所述新条件。\\
chap2:125--140 & \Ccat & Brownian/time-change 生成元系数正确，只在 $t>0,h(t)>0$；不能声称完整定义域恰为 $C_0^2$。\\
chap2:1278 & \Bcat & 换元后的分母应随新变量写为 $r-h(t)$；原式保留 $h(T)-h(t)$。\\
chap2:154--162 & \Ccat & 扩散 Copula 系数公式正确，但缺概率流边界；第一章的全局一致椭圆假设又排除 CIR。\\
chap3:55 & \Ccat & 原作者注释“这里存疑”准确指出 backward drift 问题；修正版采用 forward drift。\\
chap3:128--175 & \Ccat & Gamma 两例从一个零点跳到全域符号，是必要/充分方向错误。\\
chap3:239--241 & \Ccat & 把 $\gamma(t)<0$ 改写成常数 $\alpha>0$，却仍称回复；正确漂移是 $d(m-x)$，$d>0$。\\
chap3:706--710 & \Ccat & 均值回复命题证明把目标漂移写成 0，实际上证明了鞅而不是符号回复。\\
chap3:747--759 & \Ccat & It\^o 零漂移只直接给局部鞅；需加入 \eqref{cmp:square-int} 或 class--(D)。\\
chap3:765--791 & \Ccat & 反用充分命题，并混淆 $u$ 与物理状态；最终结论可由 \eqref{cmp:mixed-mr} 重证。\\
chap3:793--847 & \Ccat & 第 802 行重复乘事件概率；最终混合转移核仍正确。\\
chap3:849--883 & \Ccat & 生成元与 PDE 正确；Feynman--Kac 条件不完整。\\
chap3:889--906 & \Ccat & Kummer $U$ 的公式只覆盖相应非整数参数分支；整数参数需极限定义，且符号 $\Phi$ 与正态 CDF 冲突。\\
\bottomrule
\end{longtable}

\section{25 个实证对象能否嵌入修正理论}

这里的“实现”必须拆成四件事：前向映射、逆映射、转移密度和估计尝试。当前
\path{output/tables/full_model_implementation_status.csv} 中的
\texttt{density\_implemented=True} 对 SpecialFunctionTransformedOU、
SpecialFunctionTransformedCIR 和 GaussianGammaTimeVarying 写得过强：前两者只有前向/导数或有限分支扫描，后一项只有端点诊断。另需强调，注册表覆盖或保留失败诊断不等于理论有效。

本表采用五类身份：S=修正后一维严格结构；C=补条件、有限时域或改变观测对象后成立；L=潜在状态/预测扩展，不是一维结构；B=仅基准；N=以当前身份不能嵌入。

\scriptsize
\begin{longtable}{p{0.035\textwidth}p{0.205\textwidth}p{0.055\textwidth}p{0.61\textwidth}}
\toprule
序号 & Registry 名称 & 身份 & 嵌入结论、必要条件与当前实现边界\\
\midrule
\endfirsthead
\toprule
序号 & Registry 名称 & 身份 & 嵌入结论、必要条件与当前实现边界\\
\midrule
\endhead
1 & OUF & S & 标准 OU，$\kappa,\sigma>0$，使用精确 Normal 转移；满足强条件均值回复。\\
2 & ThresholdOU\_\allowbreak{}Heteroskedastic & B & 漂移方向指向阈值，但不同 $\kappa_\pm$ 时平稳均值一般不等于阈值。当前是 Euler Gaussian QMLE 而非精确转移；估计 $\kappa_+=2.0288$、$\kappa_-=0.000158$，不能满足原定义中“漂移零点=渐近均值”。可另列“方向性回复”。\\
3 & CIR\_Price & S & 理论上是标准 CIR；需非负边界条件。当前 SciPy 非中心卡方后端失败率为 100\%，实际使用条件矩 Normal fallback，故结构过程严格而经验似然不是精确 CIR 密度。\\
4 & CIR\_\allowbreak{}DepegPressure & C & 把 $D_t=|P_t-1|+\varepsilon$ 本身定义为 CIR 时成立；它不能恢复价格偏离的正负号，也不能与 signed log-price density 混为同一评分组。\\
5 & AffineTransformedOU & S & $A_0,d,\kappa>0$，全轴唯一逆与 Jacobian；直接属于扩散单调变换定理。\\
6 & Exponential\allowbreak{}TransformedOU & C & $A_0,d>0$、逐点支持；正价格还需下界非负。当前目标是 0 而非 1 美元，且估计弱识别/触边。\\
7 & SpecialFunction\allowbreak{}TransformedOU & N & 偶函数导致非单射；代码只有前向值和导数诊断，没有观测转移密度。改成平方 OU 是新模型且 Copula 改变。\\
8 & AffineTransformedCIR & S & 标准 CIR、$A_0,d>0$ 及逐时支持；理论精确，当前继承 CIR 的 100\% Normal fallback 与弱识别。\\
9 & Quadratic\allowbreak{}TransformedCIR & C & 只作为满足 \eqref{cmp:quadratic-window} 的有限时域模型；当前支持与单调违反率均为 1，不能进入全时域主定理。\\
10 & Exponential\allowbreak{}TransformedCIR & C & 满足 \eqref{cmp:cir-exp-moment}、实对数、单调及正确支持时成立；当前支持违反率约 0.94725。\\
11 & SpecialFunction\allowbreak{}TransformedCIR & C & 只有在固定有限分支并证明单调、逆映射、矩与支持后才成立。当前只扫描 16 个 $(\zeta,d)$ 组合，没有数值逆和完整 log-density。\\
12 & MOUF & S & 平稳 OU 加独立常强度 Poisson 重置，核密度精确；直接由 \eqref{cmp:mixed-mr} 嵌入。\\
13 & MCIR & S & 在完整 CIR 转移密度且无边界原子的条件下，平稳 CIR 加独立重置严格成立；当前基础 CIR 固定、100\% Normal fallback，且弱识别被显式标记。\\
14 & SeasonalIntensity\allowbreak{}MOU & S & 补非齐次 Poisson 命题后严格成立；强度还须非负、右连续且局部有界。代码用幅度半径小于 0.95 保证 $\lambda(t)>0$。当前最小强度约 0.04757/日，但参数弱识别。\\
15 & JumpOU & B & $\sigma>0$ 的精确 compound-Poisson OU 可纳入有密度的跳扩散扩展；$\sigma=0$ 的纯跳边界含“无跳”原子，不属于本文 Copula 密度范围。若跳均值为 $m_J$，回复中心应改为 $\theta+\lambda m_J/\kappa$。当前“无跳精确+至少一跳矩匹配”不是已证明的半群。\\
16 & MarkovSwitching\allowbreak{}OU & L & $(X_t,S_t)$ 是 Markov，观测 $X_t$ 一般不是。须在隐状态空间陈述；当前不同期限还使用不同层级的 HMM/endpoint mixture。\\
17 & NestedOU & L & 两 OU 因子向量 Markov；标量和除速率相同或权重退化外非 Markov。当前是逐期限二元 Gaussian composite objective，不是完整 Kalman 路径似然。\\
18 & RandomWalk & B & 明确的无回复负控制，不满足漂移符号或有限长期均值。\\
19 & GaussianGamma\allowbreak{}Stationary & N & Gaussian OU Copula 加 Gamma 边缘本身是有效平稳 Markov 密度，但尚未通过 \eqref{cmp:gamma-global}，甚至未找到有限的必要零点根；不能称为已证明均值回复模型。\\
20 & GaussianGamma\allowbreak{}TimeVarying & N & 原 ODE 只做单点诊断；当前没有求尺度路径、时变密度或估计，且收敛端点仍受 Case 1 限制。\\
21 & ParametricCopula\allowbreak{}Grid & L & 必须逐单元格分类：Gaussian $\rho(h)=e^{-\kappa h}$ 及其独立重置混合有精确半群；在容许的内部参数下，Student-$t$/Clayton/Gumbel/Frank 等只有逐期限绝对连续二元 Copula 密度，连续时间 Chapman--Kolmogorov 未证，奇异边界参数排除。离散 5 分钟核迭代可有效，但不是同参数连续时间族。\\
22 & TwoScale\allowbreak{}Gaussian & L & 两因子向量模型可严格定义；$\rho(h)=we^{-\kappa_fh}+(1-w)e^{-\kappa_sh}$ 的标量观测一般不是 Markov，不能用一维生成元定理。\\
23 & TwoScale\allowbreak{}Mixed & L & 若 Poisson 时刻重置整个二维向量，可严格构造潜在 Markov 模型；当前仅实现各期限的二元混合密度，过程级一致性需由向量构造给出。\\
24 & FullyTimeVarying\allowbreak{}MarginalCopula & C & 当确定性 $s(t)>0$ 时，$X_t=m(t)+s(t)Z_t$ 本身是一维时变 Markov 过程，Gaussian 转移密度与 Copula 密度精确。向固定 $m_\infty$ 回复还需 $\beta=\kappa-s'/s>0$、$m'=\beta(m_\infty-m)$ 及 $m(t)\to m_\infty$；当前非零趋势与永久季节项破坏的是固定目标回复，不是 Markov 性。\\
25 & MixedCopula\allowbreak{}BondPDE & S & 风险中性理论对象在补齐 \eqref{cmp:bond-pde} 的条件后严格；USDT spot 数据没有债券横截面或风险中性测度，因此本次不可估，不是理论失败。\\
\bottomrule
\end{longtable}
\normalsize

\paragraph{不能用一个“已实现”字段代替理论判定。}
建议以后对每个对象分别记录下列字段：
\begin{center}
\small\ttfamily
\begin{tabular}{@{}ll@{}}
theory\_embedding\_class & markov\_state\_dimension\\
observed\_process\_markov & chapman\_kolmogorov\_proved\\
asymptotic\_mean\_exists & global\_drift\_sign\_pass\\
transition\_density\_status & fallback\_rate\\
support\_pass & identification\_status
\end{tabular}
\end{center}
其中 \texttt{transition\_density\_status} 的取值为
exact、approximate、diagnostic 或 not implemented。

\section{最终处置清单}

\begin{enumerate}[label=\arabic*.]
  \item 保留并补条件：OU/CIR 基础条件均值、Gaussian 生成元、有限时距 Copula 恒等式、仿射 OU/CIR、指数 OU/CIR、混合核、Poisson 构造与债券 PDE。
  \item 降级为诊断：Gamma 两例的标量方程/ODE；在全域符号验证前不得写“满足均值回复”。
  \item 改成有限时域：二次 CIR；正文必须显示窗口 $T$ 与约束 \eqref{cmp:quadratic-window}。
  \item 保留为待指定族：CIR Kummer 分离解；每个分支单独审计，不能用一个模型名覆盖。
  \item 从一维单调框架删除：OU 偶函数特殊解和错误 CIR 第二变量变换。前者可另建多分支隐状态模型，后者用 \eqref{cmp:correct-cir-reduction} 替换。
  \item 将 Markov-switching、NestedOU 和两尺度模型放进潜在向量 Markov 章节；不要声称观测价格自身是一阶 Markov。
  \item 将非 Gaussian Copula grid 与近似 JumpOU 明确标为“逐期限预测密度”；除非另证半群嵌入，不进入连续时间结构定理。
  \item 风险中性债券定价和物理测度 stablecoin spot 实证分开报告。
\end{enumerate}

\paragraph{与修正版文件的关系。}
工作区中的 \path{manuscript/source/meanreversion_article_theory_revised.tex} 已按照本报告的修正原则重写为自包含理论稿：所有定理均给出可核查条件；错误推导被明确替换；实证模型以身份表而非“全部强塞进同一理论”定位。

\end{document}
